%%%%     Please only edit bits which say to edit them     %%%%
%%%% They usually start with %///% to show where they are %%%%



%%% START : DO NOT EDIT %%%%

% Document setup
\documentclass[14pt]{extarticle}

\usepackage{graphicx,geometry,hyperref,enumitem,nicefrac,extsizes,color,nameref,calc,textcomp,csquotes,setspace,titlesec}

\hypersetup{
    pdfborder={0 0 0}
}

% Important: This must come near the start of the setup file

\makeatletter
\newcommand*{\currentname}{\@currentlabelname}
\makeatother

\makeatletter
\def\normaljustify{%
  \let\\\@centercr\rightskip\z@skip \leftskip\z@skip%
  \parfillskip=0pt plus 1fil}
\makeatother

\setlength\parindent{0pt}
\setlength{\parskip}{\baselineskip}
\setcounter{secnumdepth}{2}
\setlength{\fboxsep}{0.4cm}
\setlength{\fboxrule}{0.04cm}

\definecolor{niceblue}{RGB}{166,218,255}
\definecolor{nicebluedark}{RGB}{0,104,179}
\definecolor{nicegreen}{RGB}{202,255,166}
\definecolor{nicegreendark}{RGB}{71,179,0}
\definecolor{nicered}{RGB}{255,166,166}
\definecolor{nicereddark}{RGB}{179,0,0}
\definecolor{niceyellow}{RGB}{255,254,166}
\definecolor{niceyellowdark}{RGB}{219,216,0}

\newcommand{\sectionbreak}{\clearpage}
\newcommand{\nosectionbreak}{
    \renewcommand{\sectionbreak}{%
    }
}

\newcommand{\actioncol}{niceblue}
\newcommand{\notecol}{nicered}
\newcommand{\questioncol}{niceyellow}
\newcommand{\goalcol}{nicegreen}

\newcommand{\group}[1]{
	\begin{center}
	\parbox[c]{\linewidth} {
		#1
	}
	\end{center}
}

\newcommand{\colourbox}[2]{
	\begin{center}
	\fcolorbox{#1dark}{#1}{
		\parbox[c]{.9\linewidth} {
			#2 
		}
	}
	\end{center}
}

\newcommand{\numberlist}[2][]{
	\begin{enumerate}[series=#1]
		#2
	\end{enumerate}
	\vspace{-\topsep}
}

\newcommand{\resumenumberlist}[2]{
	\begin{enumerate}[resume=#1]
		#2
	\end{enumerate}
	\vspace{-\topsep}
}

\newcommand{\itemlist}[1]{
	\begin{itemize}
		#1
	\end{itemize}
	\vspace{-\topsep}
}

\newcommand{\valuelist}[2]{
	\begin{description}[align=left,labelwidth=*,leftmargin=#1]
		#2
	\end{description}
}

\newcommand{\actions}[2][\currentname acs]{
	\colourbox{\actioncol}{ 
		{\small\color{\actioncol dark} Actions}
		\resumenumberlist{ #1 }{ #2 }
	}
}

\newcommand{\restartactions}[2][\currentname acs]{
	\colourbox{\actioncol}{ 
		{\small\color{\actioncol dark} Actions}
		\numberlist[#1]{ #2 } 
	}
}

\newcommand{\notes}[1]{
	\colourbox{\notecol}{ 
		{\small\color{\notecol dark} Notes}\par
		#1
	}
}

\newcommand{\questions}[2][questions]{
	\colourbox{\questioncol}{ 
		{\small\color{\questioncol dark} Questions}
		\resumenumberlist{ #1 }{ #2 } 
	}
}
\newcommand{\restartquestions}[2][questions]{
	\colourbox{\questioncol}{ 
		{\small\color{\questioncol dark} Questions}
		\numberlist[#1]{ #2 } 
	}
}
\newcommand{\goals}[1]{
	\colourbox{\goalcol}{ 
		{\small\color{\goalcol dark} Goals}\par
		#1 
	}
}

\newcommand{\boxesdescription}{
	\vspace{-\topsep}
	\begin{description}[align=left,labelwidth=*,leftmargin=3cm]
		\item [{\color{\actioncol dark}Actions}] Stuff for you to do. They are highlighted in blue.
		
		\item [{\color{\notecol dark}Notes}] Notes about important stuff you need to be aware of (and possibly remember!). They are highlighted in red.
		
		\item [{\color{\questioncol dark} Questions}] Questions you should try to answer. Sometimes you'll need to write things down; other times you'll need to build something in the game. They are highlighted in yellow.

		{\bfseries Ask a helper or the teacher to check your answers.}
		
		\item [{\color{\goalcol dark} Goals}] Stuff you should have completed at the end of each section. They are highlighted in green.
	\end{description}
}

\newcommand{\screenshot}[3][.9\linewidth-0.8cm]{
	\begin{center}
	\includegraphics[width=#1]{img/screenshots/#2}\par
	{\small #3}
	\end{center}
}

\newcommand{\image}[3][.9\linewidth-0.8cm]{
	\begin{center}
	\includegraphics[width=#1]{img/other/#2}\par
	{\small #3}
	\end{center}
}

\newcommand{\todo}[1]{
	\image[7cm]{todo}{#1}
}

\newcommand{\spacer}{
	\par\vspace{14pt}
}

\newcommand{\answer}[1]{

{\bfseries #1}
}

%%% END : DO NOT EDIT %%%%



%///% Workshop info - edit these values for your workshop
\def\workshoptitle{Nature By The Numbers - Intro To Data Science}
\def\workshopsubtitle{Preparation Notes for Teachers}
\def\workshopauthor{Drum Ogilvie}



%%% START : DO NOT EDIT %%%%

\begin{document}

\sffamily

% Title page - do not edit
\begin{titlepage}

	\newgeometry{margin=2cm}
	
	\begin{figure}
	\centering
	\begin{minipage}{.5\textwidth}
		\centering
		\includegraphics[width=.6\linewidth]{img/Bristol-University-Logo}
	\end{minipage}%
	\begin{minipage}{.5\textwidth}
		\centering
		\includegraphics[width=.4\linewidth]{img/Digimakers}
	\end{minipage}
	
	\vspace{18pt}
	
	\end{figure}
	
	\centering
	\normalfont
	{\Large Merchant Venturers School of Engineering\par Outreach Programme\par}
	\vspace{2cm}
	\sffamily
	{\huge\bfseries\workshoptitle\par}
	\vspace{0.5cm}
	{\LARGE\bfseries\workshopsubtitle}
	
	\vfill
	
	{\normalsize Created by\par
	\large\slshape\workshopauthor}
	
	\vfill
	
	{\normalsize Organised by\par
	\large\href{mailto:Caroline.Higgins@bristol.ac.uk}{Caroline.Higgins@bristol.ac.uk}}

	\vfill

	{\large Published on \today\par}
\end{titlepage}

% Document setup - do not edit
\newgeometry{margin=1.5cm}

\renewcommand{\abstractname}{Workshop Summary Notes}
\begin{abstract}

\vspace{-\topsep}
    
\noindent
\vspace{-\topsep}
    
\setlist{leftmargin=0mm}
\itemlist {
    \setlength{\itemsep}{18pt}
    
%%% END : DO NOT EDIT %%%%
    
% Put notes to teachers and session helpers as a bullet point list in this section

%///% Begin editable content

    \item This workshop is intended as a 90 minutes activity ??
    
    \item The workshop is intended for students age ?? . (marked as advanced)

    \item The content is intended to be learnt through following a presentation and completing exercises.

    \item The learning platforms are VirtualBox, Ansible and Vagrant.

    \item No additional software is required.

    \item The workshop is intended for students at advanced level, with at least basic understanding of Python.
		
    \item This document includes some notes about what you need to do to prepare for us to come and run the workshop.
    
% End editable content

%%% START : DO NOT EDIT %%%%

}

\end{abstract}

% Document setup - do not edit
\newgeometry{margin=1.5cm}


%%% END : DO NOT EDIT %%%%

% Sections from this point onwards are your choice, though Introduction, Conclusion (/Wrap-up/What we learnt) and Extra Resources sections are recommended.

%///% Begin editable content

\section{Guide}

The following notes are guidelines for the things we have thought of or know from experience need to be prepared before we arrive to run the workshop. Please use these notes as an addition to any guidelines and preparation you would normally follow.

\valuelist{6cm}{
    
    \item [Provision of equipment] We will provide laptops because of the required software in order to run the session. We will need to allow at least an additional 15 minutes for handing them out and logging everybody on / getting set up.
    
    \item [Software Installation] Software will be installed on the laptops provided.
    
    \item [Website Access] Please ensure your IT administrator has unblocked our Github site: https://github.com/MVSE-Outreach/ The easiest way to test this is to get a student (not a teacher) to open the websites and attempt to browse them. 
        
    \item [Students with special needs] We are happy to accommodate students who have special educational needs or disabilities. If you can, please make us aware of any students in the group who might need extra attention and any additional information that will help us make the workshop as enjoyable and successful as possible for them.
    
    \item [DBS Checks] Your school or organisation may require valid DBS checks (previously called CRB checks) for our volunteers. Please let us know if this is the case so we complete the process. All of our volunteers are current members of the University of Bristol.
    
    \item [Name List and Contact Information] We would very much like to be able to keep in touch with students to encourage them to attend our other free, open outreach events (such as Digimakers). If possible, please supply us with a list of the names of students and their parent's contact information so that we may follow up with them after the last workshop. We never send spam and we'll only contact parents to let them know about our other outreach events that are available to them.
}

\end{document}